\chapter{1994:Pierre-Louis Lions}\label{ch1994pierrelouislions:pierre-louis-lions-fields-medal-1994-a-comprehensive-report-on-his-work-and-impact}

\begin{center}\rule{0.5\linewidth}{0.5pt}\end{center}

\textbf{Prepared:} July 2026
\textbf{Subject:} Pierre-Louis Lions (born August 11, 1956, Grasse, France)
\textbf{Position:} Professor, College de France, Paris (since 2003)
\textbf{Doctoral Advisor:} Haim Brezis (Universite Pierre et Marie Curie, Paris)
\textbf{Fields Medal Citation:} \emph{``For his work on the theory of
nonlinear partial differential equations.''}

\begin{center}\rule{0.5\linewidth}{0.5pt}\end{center}

\begin{prereqs}
Partial differential equations (elliptic, parabolic, hyperbolic),
functional analysis, measure theory, optimal transport.
\end{prereqs}

\begin{objectives}
Explain viscosity solutions and their role in Hamilton--Jacobi
equations; describe the concentration--compactness principle;
outline the mean-field games framework.
\end{objectives}

% ============================================================
\section{Problem History}
% ============================================================

Nonlinear partial differential equations describe a vast range of
physical phenomena, from fluid flow to quantum mechanics to
economics.  The challenge of understanding solutions when nonlinear
terms dominate had driven research since the 19th century.

The \textbf{Hamilton--Jacobi equation} $u_t + H(Du) = 0$ arises in
optimal control and classical mechanics.  Unlike linear equations,
it generally has no classical (smooth) solutions.  The notion of
\textbf{viscosity solutions}, introduced independently by Crandall
and Lions (1983), provided the correct framework for weak solutions.

The \textbf{concentration--compactness principle} was introduced by
Lions in his doctoral thesis (1984) to handle problems where
standard compactness arguments fail due to loss of mass at infinity.
This was particularly important for problems in $\mathbb{R}^N$
where bounded sequences can escape to infinity.

The theory of \textbf{mean-field games}, developed by Lions and
Lasry starting around 2006, provides a mathematical framework for
modeling large numbers of rational agents interacting through their
aggregate effect.  This connects PDE theory, game theory, and
economics.

Lions's work on the \textbf{Boltzmann equation} (kinetic theory of
gases) and the \textbf{Navier--Stokes equations} also contributed
to the understanding of fundamental equations of mathematical
physics.

Lions's early work, done while a student of Brezis, showed
remarkable originality.  His 1982 doctoral thesis established
concentration--compactness and its applications to variational
problems, particularly the Sobolev inequality and the Yamabe
problem.

The \textbf{Yamabe problem} asks whether every Riemannian metric on
a compact manifold can be conformally rescaled to have constant
scalar curvature.  The problem reduces to solving a nonlinear
elliptic equation on the manifold.  Lions's
concentration--compactness principle provided the key compactness
argument needed to handle the non-compact setting.

% ============================================================
\section{How Was It Solved}
% ============================================================

The \textbf{viscosity solution} framework uses a comparison
principle: a continuous function $u$ is a viscosity subsolution of
$F(x, u, Du) = 0$ if for every smooth test function $\phi$
touching $u$ from above at $x_0$, we have
$F(x_0, u(x_0), D\phi(x_0)) \leq 0$.  The \textbf{comparison
principle} states that if $u \leq v$ on the boundary and both are
viscosity solutions, then $u \leq v$ everywhere.

Formally, the Crandall--Lions definition is: for a first-order
equation $F(x,u,Du)=0$, the function $u$ is a viscosity
subsolution if for any $\phi\in C^1$ with $u-\phi$ having a local
maximum at $x_0$:
\[
  F\bigl(x_0, u(x_0), D\phi(x_0)\bigr) \leq 0,
\]
and $u$ is a supersolution if $u-\phi$ has a local minimum.  This
uses only continuous functions and test functions, making it
applicable to equations with non-smooth solutions.  The key
advantage is that viscosity solutions are \textbf{unique} under
appropriate conditions, thanks to the comparison principle.

The \textbf{concentration--compactness principle} decomposes a
bounded sequence $\{u_n\}$ in $L^p(\mathbb{R}^N)$ into weakly
convergent parts and a remainder concentrating at finitely many
points:
\[
  u_n \rightharpoonup u \text{ weakly, and }
  |u_n|^p \to |u|^p + \sum_{j=1}^{J} \lambda_j\,\delta_{x_j}
\]
where $\lambda_j\geq 0$ and $\delta_{x_j}$ is a Dirac mass.  This
quantifies the loss of compactness.

For the \textbf{Boltzmann equation}
$f_t + v \cdot \nabla_x f = Q(f,f)$, Lions proved existence of
global weak solutions using a variational formulation, with the
entropy $H(f) = \int f \log f$ as the key functional analytic
tool.

The \textbf{mean-field games} framework couples a Hamilton--Jacobi
equation for optimal strategy:
\[
  u_t + H(x, Du, m) = 0
\]
with a Kolmogorov--Fokker--Planck equation for agent density:
\[
  m_t - \operatorname{div}(m\,D_p H(x,Du,m)) - \Delta m = 0.
\]
The Nash equilibrium corresponds to a coupled system of these
equations.

The concentration--compactness lemma is particularly useful in
variational problems on $\mathbb{R}^N$ where the Sobolev embedding
$H^1(\mathbb{R}^N)\hookrightarrow L^p(\mathbb{R}^N)$ is not
compact.  Lions showed that one can recover compactness by
controlling the mass that escapes to infinity, using a sharp
\textbf{Sobolev inequality} with best constant.

% ============================================================
\section{Impact of This Discovery}
% ============================================================

The viscosity solution framework revolutionized nonlinear PDE
theory.  It is now the standard tool for Hamilton--Jacobi equations,
fully nonlinear elliptic equations, and degenerate parabolic
equations.  The Crandall--Lions theory resolved many longstanding
problems.

The concentration--compactness principle is fundamental in the
calculus of variations for problems in unbounded domains, applied to
the Yamabe problem, scalar curvature, and the Schrodinger equation.

Mean-field games have become a major field at the intersection of
PDE theory, game theory, and mathematical finance, with
applications to traffic flow, economic equilibrium, systemic risk,
and crowd behavior.

Lions's work on the Boltzmann equation provided the first general
existence theorem for weak solutions, connecting PDE theory to
statistical mechanics.

Beyond specific results, Lions's ability to develop new
frameworks---viscosity solutions, concentration--compactness,
mean-field games---has set a model for mathematical innovation.

The concentration--compactness principle has been extended in
numerous directions.  It is now a standard tool in the calculus of
variations, used in problems ranging from the Sobolev inequality in
$\mathbb{R}^N$ to the study of Bose--Einstein condensation in
physics.  Its influence extends to the theory of
Gamma-convergence.

Lions's work with Carleson on the \textbf{global regularity of the
2D Euler equations} showed that vortex patches remain smooth for all
time, resolving a long-standing open problem in fluid mechanics.

% ============================================================
\section{Modern Views and Improvements}
% ============================================================

Viscosity solutions have been extended to integro-differential
equations, stochastic PDEs, and equations on manifolds.  The
comparison principle has been generalized to discontinuous solutions
and irregular coefficients.

Concentration--compactness has been refined to handle profile
decompositions and concentration at infinity, crucial for
dispersive equations.

Mean-field games have grown into a rich theory with connections to
machine learning, distributed optimization, and reinforcement
learning.  Numerical methods for mean-field games are an active
research area.

The Navier--Stokes existence and smoothness problem remains open.
Lions's tools provide frameworks for study, though complete
resolution has not been achieved.

Recent developments include the \textbf{master equation} approach
to mean-field games, introduced by Lions and Lasry, which describes
the value function of the game as a function of the measure
parameter.  This infinite-dimensional PDE provides a unified
framework for understanding Nash equilibria in mean-field games.

The \textbf{diagonal approach} to mean-field games, where one
studies the convergence of $N$-player games to the mean-field
limit, has been developed by Carmona and Delarue.  This rigorous
passage from finite to infinite agent models connects the theory to
stochastic control and game theory.

Applications of mean-field games to \textbf{climate change
economics}, \textbf{epidemic modeling}, and \textbf{crowd dynamics}
have demonstrated the practical impact of Lions's theoretical
framework.  The COVID-19 pandemic inspired new applications of
mean-field game models to understand the interplay between
individual behavior and epidemic spread.

Modern work continues to develop mean-field games with common
noise, mean-field control problems, and connections to optimal
transport.

% ============================================================
\begin{center}\rule{0.5\linewidth}{0.5pt}\end{center}

\section*{Summary}
\addcontentsline{toc}{section}{Summary}

Lions developed viscosity solutions for nonlinear PDEs, the
concentration--compactness principle, and the mean-field games
framework.  These three innovations fundamentally shaped modern
analysis and its applications to physics, economics, and
probability.

% ============================================================
